\documentclass[12pt]{article}
\usepackage{amsmath, amssymb, amsthm}
\usepackage{array}
\usepackage{booktabs}
\usepackage[margin=1in]{geometry}
\setlength{\headheight}{14.5pt}
\addtolength{\topmargin}{-2.5pt}
\usepackage{xcolor}
\usepackage{tcolorbox}
\tcbuselibrary{breakable}
\usepackage{enumitem}
\usepackage{fancyhdr}
\usepackage{graphicx}

% Define solution environment
\newtcolorbox{solution}{
    colback=blue!5!white,
    colframe=blue!75!black,
    title=Solution,
    fonttitle=\bfseries,
    breakable
}

\newtcolorbox{explanation}{
    colback=green!5!white,
    colframe=green!75!black,
    title=Explanation,
    fonttitle=\bfseries,
    breakable
}

\DeclareMathOperator{\Tr}{Tr}
\DeclareMathOperator{\Var}{Var}
\DeclareMathOperator{\Cov}{Cov}

\title{CHAPTER-WISE PROBLEMS\\
\large Markov Chains, Monte Carlo \& Deterministic Models\\
\small MIT/Stanford Level}
\author{Statistical Foundation of Data Science\\
CSU1658}
\date{December 2025}

\pagestyle{fancy}
\fancyhf{}
\rhead{Stochastic Models}
\lhead{CSU1658}
\cfoot{\thepage}

\begin{document}

\maketitle

\tableofcontents
\newpage

% ============================================================
% MARKOV CHAIN PROBLEMS
% ============================================================

\section{Markov Chains}

\subsection*{Problem 1: Credit Rating Migration Model}

A credit rating agency models corporate bond ratings using a discrete-time Markov chain with 5 states: AAA (1), AA (2), A (3), BBB (4), Default (5). Historical data from 500 companies over one year yields the following transition matrix:

\begin{center}
\begin{tabular}{cccccc}
\toprule
From/To & AAA & AA & A & BBB & Default \\
\midrule
AAA & 0.92 & 0.06 & 0.01 & 0.01 & 0.00 \\
AA & 0.03 & 0.88 & 0.08 & 0.01 & 0.00 \\
A & 0.00 & 0.05 & 0.85 & 0.08 & 0.02 \\
BBB & 0.00 & 0.01 & 0.08 & 0.82 & 0.09 \\
Default & 0.00 & 0.00 & 0.00 & 0.00 & 1.00 \\
\bottomrule
\end{tabular}
\end{center}

Current portfolio distribution: 150 AAA bonds, 200 AA bonds, 100 A bonds, 50 BBB bonds.

\begin{enumerate}[label=(\alph*)]
\item Calculate the expected distribution of ratings after 1 year by multiplying the initial distribution vector by the transition matrix. How many bonds are expected to default in the first year?

\item Compute the 2-step transition probabilities by squaring the transition matrix. What is the probability that an AAA-rated bond will be BBB-rated or worse after 2 years? Compare this to the 1-year probability.

\item Identify transient and absorbing states. Prove that Default is an absorbing state by showing $P(\text{Default} \to \text{Default}) = 1$. Are there any other absorbing states?

\item Calculate the expected time until default for a bond currently rated A using first-step analysis. Set up the system of equations: $h_i = 1 + \sum_{j \neq 5} P_{ij}h_j$ where $h_5 = 0$ (Default state). Solve for $h_3$ (state A).

\item The portfolio manager wants to know the long-term default probability starting from each non-default state. Calculate absorption probabilities $f_i$ (probability of eventual default from state $i$) by solving: $f_i = P_{i5} + \sum_{j=1}^{4} P_{ij}f_j$ with $f_5 = 1$. Order states by riskiness.

\item Time-value adjustment: If the annual discount rate is 5\%, calculate the expected present value of \$100 bond cash flows assuming annual coupon payments until default or maturity (10 years). For an A-rated bond, compute: $E[PV] = \sum_{t=1}^{10} (1.05)^{-t} \times P(\text{survive to year } t) \times \$8 + (1.05)^{-10} \times P(\text{no default}) \times \$100$.
\end{enumerate}

\subsection*{Problem 2: Web Page Ranking with PageRank Algorithm}

A simplified web network has 6 pages with the following link structure:

\begin{center}
\begin{tabular}{ccl}
\toprule
Page & Outgoing Links & Destinations \\
\midrule
1 & 2 & Pages 2, 3 \\
2 & 3 & Pages 1, 4, 5 \\
3 & 2 & Pages 2, 6 \\
4 & 1 & Page 5 \\
5 & 3 & Pages 1, 2, 6 \\
6 & 2 & Pages 3, 5 \\
\bottomrule
\end{tabular}
\end{center}

Transition probabilities: $P_{ij} = 1/(\text{outdegree of page } i)$ if page $i$ links to page $j$, else 0.

Random surfer model with damping factor $d = 0.85$: At each step, with probability 0.85 follow a random outlink, with probability 0.15 jump to any page uniformly at random.

\begin{enumerate}[label=(\alph*)]
\item Construct the 6×6 transition matrix $\mathbf{P}$ for the pure random walk (without damping). Verify that each row sums to 1. Identify any dangling nodes (pages with no outlinks).

\item Apply the damping factor to get the Google matrix: $\mathbf{G} = 0.85\mathbf{P} + 0.15\mathbf{E}$ where $\mathbf{E}$ is the $6 \times 6$ matrix with all entries equal to $1/6$. Calculate the first row of $\mathbf{G}$.

\item Find the stationary distribution $\boldsymbol{\pi}$ satisfying $\boldsymbol{\pi}^T = \boldsymbol{\pi}^T\mathbf{G}$ and $\sum \pi_i = 1$. Use power iteration: start with uniform distribution $\boldsymbol{\pi}^{(0)} = (1/6, \ldots, 1/6)$ and iterate $\boldsymbol{\pi}^{(k+1)} = \mathbf{G}^T\boldsymbol{\pi}^{(k)}$ until convergence (change < 0.001). Report the PageRank scores.

\item Rank pages by their stationary probabilities. Which page has highest PageRank? Calculate the ratio of highest to lowest PageRank. How does this relate to the link structure?

\item Sensitivity analysis: Page 2 adds a new link to Page 6. Recalculate the transition matrix and stationary distribution. How much does Page 6's PageRank increase? What about Page 2's rank?

\item Convergence rate: The mixing time is related to the second-largest eigenvalue magnitude $|\lambda_2|$ of $\mathbf{G}$. Given $|\lambda_2| \approx 0.85$ (the damping factor), estimate the number of iterations needed for $\|\boldsymbol{\pi}^{(k)} - \boldsymbol{\pi}\| < 0.001$ using the bound: iterations $\approx \frac{\ln(6 \times 0.001)}{\ln(0.85)}$.
\end{enumerate}

\subsection*{Problem 3: Queue Management at Service Center}

A customer service center has 3 agents. Customers arrive and form a single queue. The system state is the number of customers in the system (0, 1, 2, 3, 4, 5+, capped at 5). Transition probabilities per minute:

\begin{center}
\begin{tabular}{ccccccc}
\toprule
From $\backslash$ To & 0 & 1 & 2 & 3 & 4 & 5 \\
\midrule
0 & 0.70 & 0.25 & 0.04 & 0.01 & 0.00 & 0.00 \\
1 & 0.50 & 0.30 & 0.15 & 0.04 & 0.01 & 0.00 \\
2 & 0.30 & 0.35 & 0.20 & 0.10 & 0.04 & 0.01 \\
3 & 0.25 & 0.30 & 0.25 & 0.12 & 0.06 & 0.02 \\
4 & 0.20 & 0.25 & 0.25 & 0.15 & 0.10 & 0.05 \\
5 & 0.15 & 0.20 & 0.25 & 0.20 & 0.12 & 0.08 \\
\bottomrule
\end{tabular}
\end{center}

Performance metrics from simulation: Average time between arrivals: 2.4 minutes. Average service time: 6.5 minutes.

\begin{enumerate}[label=(\alph*)]
\item Determine if this chain is irreducible by checking if all states communicate. Can you reach state 5 from state 0? Can you return to state 0 from state 5? Construct a communication diagram.

\item Find the stationary distribution $\boldsymbol{\pi}$ by solving $\boldsymbol{\pi}\mathbf{P} = \boldsymbol{\pi}$ subject to $\sum \pi_i = 1$. This is equivalent to solving the system: $\pi_j = \sum_{i=0}^{5} \pi_i P_{ij}$ for $j=0,\ldots,5$ plus the normalization constraint. Use the iterative method or solve the linear system.

\item Calculate the expected number of customers in the system: $L = \sum_{i=0}^{5} i \times \pi_i$. Using Little's Law $L = \lambda W$ where $\lambda = 1/2.4$ customers per minute, compute the expected waiting time $W$ in the system.

\item Compute the probability that an arriving customer faces a full queue (state 5) and must wait or be turned away: $\pi_5$. If the center operates 8 hours per day (480 minutes), how many customers per day experience delays?

\item Capacity planning: Management considers adding a 4th agent, which changes transition probabilities (service becomes faster). New stationary distribution: $\boldsymbol{\pi}' = (0.75, 0.18, 0.05, 0.015, 0.004, 0.001)$. Calculate the percentage reduction in expected queue length. Is the improvement worth the cost?

\item Variance of queue length: Calculate $\text{Var}(N) = \sum_{i=0}^{5} i^2\pi_i - L^2$ for both 3-agent and 4-agent configurations. How does increased capacity affect variability? Use the coefficient of variation $CV = \sqrt{\text{Var}(N)}/L$ to compare relative variability.
\end{enumerate}

% ============================================================
% MONTE CARLO PROBLEMS
% ============================================================

\section{Monte Carlo Methods}

\subsection*{Problem 1: Option Pricing via Monte Carlo Simulation}

A European call option on a stock has strike price $K = \$50$ and expires in 1 year ($T = 1$). Current stock price $S_0 = \$48$. The stock follows geometric Brownian motion: $S_T = S_0 \exp[(r - \sigma^2/2)T + \sigma\sqrt{T}Z]$ where $Z \sim N(0,1)$, risk-free rate $r = 0.05$, volatility $\sigma = 0.30$.

Monte Carlo simulation with $N = 10,000$ paths yields:

\begin{center}
\begin{tabular}{ccccc}
\toprule
Statistic & Min & Q1 & Median & Q3 & Max \\
\midrule
Final Price $S_T$ & 18.42 & 38.65 & 49.28 & 62.14 & 125.67 \\
Payoff $(S_T - K)^+$ & 0.00 & 0.00 & 0.00 & 12.14 & 75.67 \\
\bottomrule
\end{tabular}
\end{center}

Sample mean payoff: $\bar{X} = 6.82$, Sample standard deviation: $s = 10.45$

\begin{enumerate}[label=(\alph*)]
\item Estimate the option value by discounting the average payoff: $V_0 = e^{-rT} \times \bar{X}$. Calculate the present value using $r = 0.05$ and $T = 1$.

\item Construct a 95\% confidence interval for the true option value using $V_0 \pm 1.96 \times \frac{s}{\sqrt{N}} \times e^{-rT}$. Calculate the margin of error. What does this interval tell us about pricing accuracy?

\item Variance reduction via antithetic variates: For each random $Z_i$, also simulate with $-Z_i$. This produces pairs $(S_T^i, S_T^{-i})$ with negative correlation. If the antithetic method reduces variance by 40\%, calculate the effective sample size: how many standard simulations would be needed to achieve the same precision as 10,000 antithetic pairs?

\item Calculate the probability that the option expires in-the-money (payoff > 0) as the fraction of paths where $S_T > K$. From the simulation, approximately what percentage of paths resulted in positive payoff? Using this, estimate the probability $P(S_T > 50)$.

\item Greeks estimation: Delta measures option sensitivity to stock price changes. Using finite difference: $\Delta \approx \frac{V(S_0 + h) - V(S_0 - h)}{2h}$ with $h = 0.01$. If simulations at $S_0 = 48.01$ and $S_0 = 47.99$ give values 6.84 and 6.80 respectively, calculate delta. Interpret: for a \$1 increase in stock price, how much does option value change?

\item Compare to Black-Scholes formula value: $V_{BS} = 6.78$ (using analytical formula). Calculate the relative error: $|V_0 - V_{BS}|/V_{BS}$. Is the Monte Carlo estimate within the 95\% confidence interval? Discuss sources of simulation error (discretization, random sampling, finite paths).
\end{enumerate}

\subsection*{Problem 2: Risk Assessment via Monte Carlo Integration}

A portfolio contains investments in 3 correlated assets with returns modeled as multivariate normal. Expected annual returns: $\boldsymbol{\mu} = (0.08, 0.12, 0.06)$. Covariance matrix:

\begin{center}
\begin{tabular}{cccc}
\toprule
& Asset 1 & Asset 2 & Asset 3 \\
\midrule
Asset 1 & 0.0400 & 0.0180 & 0.0120 \\
Asset 2 & 0.0180 & 0.0900 & 0.0090 \\
Asset 3 & 0.0120 & 0.0090 & 0.0225 \\
\bottomrule
\end{tabular}
\end{center}

Portfolio weights: $\mathbf{w} = (0.4, 0.3, 0.3)$. Initial investment: \$1,000,000.

Monte Carlo simulation ($N = 50,000$ scenarios) of annual returns:

\begin{center}
\begin{tabular}{cccc}
\toprule
Percentile & Portfolio Return & Portfolio Value & Loss \\
\midrule
1\% & -0.198 & \$802,000 & \$198,000 \\
5\% & -0.092 & \$908,000 & \$92,000 \\
50\% & 0.084 & \$1,084,000 & - \\
95\% & 0.268 & \$1,268,000 & - \\
99\% & 0.352 & \$1,352,000 & - \\
\bottomrule
\end{tabular}
\end{center}

\begin{enumerate}[label=(\alph*)]
\item Calculate 95\% Value at Risk (VaR): the 5th percentile loss. From the simulation output, what is the maximum loss expected with 95\% confidence? Express as both dollar amount and percentage of initial investment.

\item Compute 95\% Conditional Value at Risk (CVaR): the expected loss given that loss exceeds VaR. Average all simulated losses that are worse than the 5th percentile. If the mean of the worst 5\% scenarios is a loss of \$124,000, calculate CVaR and compare to VaR.

\item Estimate the probability of losing more than 15\% (portfolio value below \$850,000) as the fraction of simulated scenarios with returns below -0.15. If 3,200 out of 50,000 scenarios resulted in such losses, calculate this tail probability.

\item Component VaR: Decompose total VaR into contributions from each asset. For Asset 1, calculate marginal VaR: $\text{MVaR}_1 = \frac{\partial \text{VaR}}{\partial w_1}$. If increasing $w_1$ by 0.01 (1\%) increases VaR by \$850, compute marginal VaR. The component VaR is: $\text{CVaR}_1 = w_1 \times \text{MVaR}_1$. Which asset contributes most to portfolio risk?

\item Stress testing: A market crash scenario triples the volatilities (multiply covariance matrix by 9) while maintaining correlations. Re-run the simulation with adjusted parameters. If new 5\% VaR becomes \$276,000, calculate the VaR multiplier under stress. How does portfolio risk scale with volatility?

\item Importance sampling: To better estimate tail probabilities, shift the mean return down by one standard deviation: $\boldsymbol{\mu}' = \boldsymbol{\mu} - \boldsymbol{\sigma}_p$ where $\boldsymbol{\sigma}_p$ is the portfolio volatility. This generates more extreme scenarios. If importance sampling estimates $P(\text{loss} > 15\%) = 0.068$ compared to crude Monte Carlo's 0.064, calculate the relative efficiency: $\text{Efficiency} = \frac{s^2}{s'^2}$ where $s'^2$ is the importance sampling variance.
\end{enumerate}

\subsection*{Problem 3: Reliability Analysis via Monte Carlo Simulation}

A manufacturing system has 5 components arranged in series-parallel configuration:
- Subsystem A: Components 1 and 2 in parallel (system works if at least one works)
- Subsystem B: Components 3, 4, 5 in series (all must work)
- Overall: A and B in series

Component reliabilities (probability of functioning for 1000 hours):

\begin{center}
\begin{tabular}{cccc}
\toprule
Component & Reliability & Mean Lifetime (hrs) & Std Dev (hrs) \\
\midrule
1 & 0.92 & 1200 & 180 \\
2 & 0.88 & 1100 & 200 \\
3 & 0.95 & 1500 & 150 \\
4 & 0.90 & 1300 & 220 \\
5 & 0.93 & 1400 & 190 \\
\bottomrule
\end{tabular}
\end{center}

Monte Carlo simulation ($N = 100,000$ trials) results:

\begin{center}
\begin{tabular}{ccc}
\toprule
Configuration & Estimated Reliability & 95\% CI \\
\midrule
Subsystem A (1,2 parallel) & 0.9896 & [0.9886, 0.9906] \\
Subsystem B (3,4,5 series) & 0.7979 & [0.7956, 0.8002] \\
Overall system & 0.7896 & [0.7873, 0.7919] \\
\bottomrule
\end{tabular}
\end{center}

\begin{enumerate}[label=(\alph*)]
\item Verify the parallel subsystem reliability using the formula: $R_A = 1 - (1-r_1)(1-r_2)$ where $r_1 = 0.92$, $r_2 = 0.88$. Compare to simulation result 0.9896. Calculate the absolute error and relative error.

\item Calculate series subsystem reliability analytically: $R_B = r_3 \times r_4 \times r_5 = 0.95 \times 0.90 \times 0.93$. Compare to simulation result 0.7979. Why does parallel configuration provide redundancy while series creates vulnerability?

\item Estimate overall system reliability and compare to the analytical solution: $R_{\text{system}} = R_A \times R_B$. Given simulation result 0.7896, is this consistent with the theoretical value? The 95\% CI width indicates precision—calculate margin of error.

\item Component importance: Which component has greatest impact on system reliability? Calculate Birnbaum importance: $I_i = \frac{\partial R_{\text{system}}}{\partial r_i}$. For component 1 in parallel: $I_1 = R_B \times (1-r_2) = 0.7979 \times 0.12$. For component 3 in series: $I_3 = R_A \times r_4 \times r_5$. Compare these sensitivities.

\item Mean Time To Failure (MTTF): Simulate component lifetimes from normal distributions (truncated at 0). For each trial, record system failure time = minimum of (Subsystem A failure time, Subsystem B failure time). If simulation yields mean system lifetime 1050 hours with standard deviation 210 hours, construct a 95\% confidence interval for true MTTF.

\item Cost-benefit analysis: Improving component 4's reliability from 0.90 to 0.95 costs \$50,000. System downtime costs \$10,000 per hour. Calculate expected annual savings (8760 hours per year): Annual cost without improvement vs. with improvement. If new system reliability becomes 0.8318, compute reduction in expected downtime hours and determine payback period.
\end{enumerate}

% ============================================================
% DETERMINISTIC MODEL PROBLEMS
% ============================================================

\section{Deterministic Models}

\subsection*{Problem 1: Epidemic Spread Model (SIR)}

A city with population 50,000 experiences a flu outbreak. The SIR (Susceptible-Infected-Recovered) model divides the population into three groups with dynamics:

Rate of new infections: $\beta = 0.00004$ (per day per contact)

Rate of recovery: $\gamma = 0.1$ (per day)

Basic reproduction number: $R_0 = \beta N / \gamma$ where $N = 50000$

Initial conditions (Day 0): $S_0 = 49,900$, $I_0 = 100$, $R_0 = 0$

Simulation results at key timepoints:

\begin{center}
\begin{tabular}{ccccc}
\toprule
Day & Susceptible & Infected & Recovered & New Cases \\
\midrule
0 & 49,900 & 100 & 0 & - \\
10 & 47,250 & 1,850 & 900 & 265 \\
20 & 38,100 & 6,200 & 5,700 & 892 \\
30 & 24,300 & 9,800 & 15,900 & 1,368 \\
40 & 12,400 & 8,600 & 29,000 & 1,220 \\
50 & 5,800 & 4,200 & 40,000 & 750 \\
60 & 2,900 & 1,100 & 46,000 & 290 \\
\bottomrule
\end{tabular}
\end{center}

\begin{enumerate}[label=(\alph*)]
\item Calculate the basic reproduction number $R_0 = \beta N / \gamma = 0.00004 \times 50000 / 0.1$. Interpret: on average, how many new infections does one infected person cause in a fully susceptible population? Is $R_0 > 1$ indicating epidemic growth?

\item Determine the peak infection day by finding when $dI/dt = 0$. From the differential equations: $dI/dt = \beta SI - \gamma I$. At the peak, $\beta SI = \gamma I$, so $S = \gamma/\beta = 0.1/0.00004$. Using the simulation data, estimate which day had maximum infected individuals. What was the peak prevalence (percentage of population infected)?

\item Calculate cumulative attack rate: the total proportion of population that eventually gets infected. From day 60 data, this is $(I_{60} + R_{60})/N = (1100 + 46000)/50000$. What percentage of the population avoided infection entirely?

\item Herd immunity threshold: Calculate the critical proportion that must be immune to stop epidemic spread: $p_c = 1 - 1/R_0$. Compare this theoretical threshold to the final susceptible proportion $S_{\infty}/N$. How close is the simulation result to the theoretical prediction?

\item Intervention analysis: A vaccination campaign at day 15 immunizes 10,000 susceptible individuals (they move to Recovered). Adjust $S_{15}$ and $R_{15}$ accordingly. Predict the new peak infection count using the adjusted initial conditions for day 15 onward. Estimate the number of infections prevented by this intervention.

\item Time-to-peak sensitivity: How does $R_0$ affect epidemic timeline? If transmission rate $\beta$ decreases by 20\% (social distancing), new $R_0' = 0.8 \times R_0$. Use the approximation: peak time $t_{\text{peak}} \propto \frac{1}{\gamma(R_0 - 1)}$. Calculate the change in peak time and interpret the public health implications of intervention delays.
\end{enumerate}

\subsection*{Problem 2: Population Growth Model with Harvesting}

A fish population in a lake follows logistic growth with harvesting: $\frac{dN}{dt} = rN(1 - N/K) - H$ where $r$ is intrinsic growth rate, $K$ is carrying capacity, and $H$ is harvest rate.

Parameters: $r = 0.4$ per year, $K = 100,000$ fish, current population $N_0 = 60,000$.

Harvest scenarios:

\begin{center}
\begin{tabular}{cccc}
\toprule
Scenario & Harvest Rate $H$ & Equilibrium $N^*$ & Time to Equilibrium \\
\midrule
No Harvest & 0 & 100,000 & 8.5 years \\
Sustainable & 10,000/yr & 75,000 & 5.2 years \\
Moderate & 18,000/yr & 55,000 & 4.8 years \\
Heavy & 25,000/yr & 0 (collapse) & 12 years \\
Critical & 10,000/yr & 50,000 & 6.1 years \\
\bottomrule
\end{tabular}
\end{center}

\begin{enumerate}[label=(\alph*)]
\item Find equilibrium points by solving $rN(1 - N/K) = H$. For sustainable harvest $H = 10,000$, solve: $0.4N(1 - N/100000) = 10000$. This gives quadratic: $0.4N - 0.000004N^2 = 10000$. Calculate both equilibrium points and determine which is stable.

\item Maximum Sustainable Yield (MSY): The maximum harvest rate that maintains equilibrium occurs at $N^* = K/2$. Calculate $H_{\text{MSY}} = rK/4 = 0.4 \times 100000 / 4$. Compare to the harvest scenarios in the table. Which scenarios exceed MSY and risk population collapse?

\item Stability analysis: For equilibrium $N^*$, calculate $\frac{d}{dN}[rN(1 - N/K) - H]$ evaluated at $N^*$. If this derivative is negative, equilibrium is stable. For $N^* = 75,000$ under sustainable harvesting, compute: $r(1 - 2N^*/K) = 0.4(1 - 150000/100000)$. Is this equilibrium stable?

\item Recovery time: Starting from depleted population $N_0 = 30,000$ with no harvesting ($H = 0$), estimate time to reach 90\% of carrying capacity. Use the logistic growth solution: $N(t) = \frac{K N_0}{N_0 + (K - N_0)e^{-rt}}$. Solve for $t$ when $N(t) = 0.9K = 90,000$.

\item Economic optimization: Each fish sells for \$15. Annual discount rate is 8\%. Calculate present value of perpetual sustainable yield: $PV = \frac{H \times \text{price}}{r_{\text{discount}}} = \frac{10000 \times 15}{0.08}$. Compare this to the one-time value of harvesting entire population: $100000 \times 15$. Which strategy maximizes long-term value?

\item Stochastic effects: Environmental variability causes growth rate to fluctuate: $r \sim \text{Uniform}(0.3, 0.5)$ per year. Using Monte Carlo simulation over 20 years with $H = 10,000$ per year, if 15\% of simulations result in population below 20,000 (threatened status), calculate the probability of population decline. Should harvest rate be reduced for precautionary reasons?
\end{enumerate}

\subsection*{Problem 3: Chemical Reactor Dynamics}

A continuous stirred-tank reactor (CSTR) processes chemical A into product B via first-order reaction: $A \to B$ with rate constant $k = 0.5$ per minute. The reactor has volume $V = 100$ liters with inflow rate $Q = 10$ L/min.

Mass balance for concentration of A: $\frac{dC_A}{dt} = \frac{Q}{V}(C_{A,\text{in}} - C_A) - kC_A$

Input concentration: $C_{A,\text{in}} = 2.0$ mol/L

Steady-state and transient data:

\begin{center}
\begin{tabular}{cccc}
\toprule
Time (min) & $C_A$ (mol/L) & $C_B$ (mol/L) & Conversion (\%) \\
\midrule
0 & 0.00 & 0.00 & 0 \\
2 & 0.52 & 0.38 & 19 \\
5 & 0.92 & 0.87 & 43.5 \\
10 & 1.18 & 1.48 & 59 \\
20 & 1.29 & 1.91 & 64.5 \\
40 & 1.32 & 2.05 & 66 \\
$\infty$ & 1.33 & 2.07 & 66.5 \\
\bottomrule
\end{tabular}
\end{center}

\begin{enumerate}[label=(\alph*)]
\item Calculate steady-state concentration of A by setting $\frac{dC_A}{dt} = 0$: solve $\frac{Q}{V}(C_{A,\text{in}} - C_A^*) = kC_A^*$ for $C_A^*$. Using $Q/V = 0.1$ per minute and $k = 0.5$ per minute, compute $C_A^* = \frac{C_{A,\text{in}}}{1 + kV/Q} = \frac{2.0}{1 + 5}$. Compare to simulation result 1.33 mol/L.

\item Determine conversion at steady state: $X = \frac{C_{A,\text{in}} - C_A^*}{C_{A,\text{in}}} \times 100\%$. From calculated $C_A^* = 1.33$, conversion is $(2.0 - 1.33)/2.0 = 0.335$ or 33.5\%. Note discrepancy with table value 66.5\%—reconcile by considering product formation: true conversion uses $C_B/C_{A,\text{in}}$.

\item Calculate time constant for reactor approach to steady state. The solution to the differential equation is: $C_A(t) = C_A^* + (C_A(0) - C_A^*)\exp(-t/\tau)$ where time constant $\tau = 1/(k + Q/V) = 1/(0.5 + 0.1) = 1/0.6$ minutes. After how many time constants (multiples of $\tau$) does the reactor reach 95\% of steady state?

\item Residence time distribution: Average time a molecule spends in reactor is $\tau_{\text{res}} = V/Q = 100/10 = 10$ minutes. The fraction of molecules with residence time less than $t$ follows: $F(t) = 1 - e^{-t/\tau_{\text{res}}}$. Calculate the percentage of molecules that exit within 5 minutes (half the average residence time).

\item Damköhler number: This dimensionless parameter compares reaction rate to flow rate: $Da = \frac{kV}{Q} = \frac{0.5 \times 100}{10} = 5$. Interpret: $Da >> 1$ means reaction-limited (fast flow, slow reaction), $Da << 1$ means mixing-limited. Is this reactor reaction-limited or mixing-limited? What does this imply for conversion efficiency?

\item Process optimization: To increase conversion to 80\%, either decrease flow rate $Q$ or increase reaction rate $k$ (higher temperature). If doubling temperature increases $k$ to 1.0 per minute (Arrhenius law), calculate new steady-state conversion using $C_A^* = C_{A,\text{in}}/(1 + kV/Q)$. Alternatively, if $Q$ is reduced to 5 L/min (keeping $k = 0.5$), compute the new conversion. Which approach is more effective?
\end{enumerate}

\end{document}
